% ============================================================
% Question Bank: ch04-02 Simplifying Expressions by Combining Like Terms
% Topic Title: Simplifying Expressions by Combining Like Terms
% CCSS: 7.EE.A.1
% Questions: 27 (15 MC + 6 SA + 6 Graphical)
% ============================================================

% --- Multiple Choice Questions (15) ---

\begin{practiceQuestion}{4-2-q01}{mc}
\begin{questionText}
Which pair of terms are like terms?
\end{questionText}
\choiceA{$6a$ and $6b$}
\choiceB{$-3x^2$ and $5x$}
\choiceC{$9y$ and $-y$}
\choiceD{$2$ and $2n$}
\correctAnswer{C}
\explanation{Like terms have the same variable raised to the same power. $9y$ and $-y$ both have the variable $y$ to the first power, so they are like terms.}
\end{practiceQuestion}

\begin{practiceQuestion}{4-2-q02}{mc}
\begin{questionText}
Simplify $4n + 9n - 5$.
\end{questionText}
\choiceA{$13n - 5$}
\choiceB{$8n$}
\choiceC{$13n + 5$}
\choiceD{$18n$}
\correctAnswer{A}
\explanation{Combine like variable terms: $4n + 9n = 13n$. The constant $-5$ has no like term. Result: $13n - 5$.}
\end{practiceQuestion}

\begin{practiceQuestion}{4-2-q03}{mc}
\begin{questionText}
Simplify $7a - 2 + 3a - 8$.
\end{questionText}
\choiceA{$10a - 10$}
\choiceB{$4a - 10$}
\choiceC{$10a + 10$}
\choiceD{$10a - 6$}
\correctAnswer{A}
\explanation{Combine variable terms: $7a + 3a = 10a$. Combine constants: $-2 - 8 = -10$. Result: $10a - 10$.}
\end{practiceQuestion}

\begin{practiceQuestion}{4-2-q04}{mc}
\begin{questionText}
Simplify $-5x + 12 + 8x - 7$.
\end{questionText}
\choiceA{$-3x + 5$}
\choiceB{$3x + 5$}
\choiceC{$13x + 5$}
\choiceD{$3x + 19$}
\correctAnswer{B}
\explanation{Combine variable terms: $-5x + 8x = 3x$. Combine constants: $12 - 7 = 5$. Result: $3x + 5$.}
\end{practiceQuestion}

\begin{practiceQuestion}{4-2-q05}{mc}
\begin{questionText}
What is the coefficient of $p$ in the expression $9 - 4p + 1$?
\end{questionText}
\choiceA{$9$}
\choiceB{$4$}
\choiceC{$-4$}
\choiceD{$1$}
\correctAnswer{C}
\explanation{The coefficient is the number multiplied by the variable. In $-4p$, the coefficient is $-4$.}
\end{practiceQuestion}

\begin{practiceQuestion}{4-2-q06}{mc}
\begin{questionText}
Simplify $3m + 2 - 3m + 6$.
\end{questionText}
\choiceA{$6m + 8$}
\choiceB{$8$}
\choiceC{$0$}
\choiceD{$-6m + 8$}
\correctAnswer{B}
\explanation{Combine variable terms: $3m - 3m = 0$. Combine constants: $2 + 6 = 8$. Result: $8$.}
\end{practiceQuestion}

\begin{practiceQuestion}{4-2-q07}{mc}
\begin{questionText}
Which expression is equivalent to $-x + 5 + 6x - 9$?
\end{questionText}
\choiceA{$5x - 4$}
\choiceB{$7x - 4$}
\choiceC{$-5x - 4$}
\choiceD{$5x + 14$}
\correctAnswer{A}
\explanation{Combine variable terms: $-x + 6x = 5x$. Combine constants: $5 - 9 = -4$. Result: $5x - 4$.}
\end{practiceQuestion}

\begin{practiceQuestion}{4-2-q08}{mc}
\begin{questionText}
Simplify $\frac{3}{5}y + \frac{2}{5}y - 1$.
\end{questionText}
\choiceA{$y - 1$}
\choiceB{$\frac{5}{10}y - 1$}
\choiceC{$\frac{1}{5}y - 1$}
\choiceD{$\frac{6}{25}y - 1$}
\correctAnswer{A}
\explanation{Add the coefficients: $\frac{3}{5} + \frac{2}{5} = \frac{5}{5} = 1$. Result: $y - 1$.}
\end{practiceQuestion}

\begin{practiceQuestion}{4-2-q09}{mc}
\begin{questionText}
A student simplifies $6x - 4 + 2x$ and writes $4x - 4$. What mistake did the student make?
\end{questionText}
\choiceA{Subtracted $2x$ from $6x$ instead of adding}
\choiceB{Forgot the constant term}
\choiceC{Multiplied the coefficients instead of adding}
\choiceD{Combined unlike terms}
\correctAnswer{A}
\explanation{The student computed $6x - 2x = 4x$ instead of $6x + 2x = 8x$. The correct answer is $8x - 4$.}
\end{practiceQuestion}

\begin{practiceQuestion}{4-2-q10}{mc}
\begin{questionText}
Simplify $1.5w - 3.2 + 2.5w + 0.8$.
\end{questionText}
\choiceA{$4w + 4$}
\choiceB{$4w - 2.4$}
\choiceC{$-1w - 4$}
\choiceD{$4w - 4$}
\correctAnswer{B}
\explanation{Combine variable terms: $1.5w + 2.5w = 4w$. Combine constants: $-3.2 + 0.8 = -2.4$. Result: $4w - 2.4$.}
\end{practiceQuestion}

\begin{practiceQuestion}{4-2-q11}{mc}
\begin{questionText}
How many terms are in the expression $4x + 7 - 2x + 3$ after simplifying?
\end{questionText}
\choiceA{$4$}
\choiceB{$3$}
\choiceC{$2$}
\choiceD{$1$}
\correctAnswer{C}
\explanation{Simplify: $4x - 2x = 2x$ and $7 + 3 = 10$. Result: $2x + 10$, which has $2$ terms.}
\end{practiceQuestion}

\begin{practiceQuestion}{4-2-q12}{mc}
\begin{questionText}
Simplify $-8k + 3k + 2k$.
\end{questionText}
\choiceA{$-3k$}
\choiceB{$3k$}
\choiceC{$-13k$}
\choiceD{$13k$}
\correctAnswer{A}
\explanation{$-8k + 3k = -5k$, then $-5k + 2k = -3k$.}
\end{practiceQuestion}

\begin{practiceQuestion}{4-2-q13}{mc}
\begin{questionText}
Which expression is already in simplest form?
\end{questionText}
\choiceA{$5a - 3a + 1$}
\choiceB{$7x - 2 + 4$}
\choiceC{$6n + 5$}
\choiceD{$12 - 4 + 3y$}
\correctAnswer{C}
\explanation{$6n + 5$ has no like terms to combine. All other choices have terms that can still be simplified.}
\end{practiceQuestion}

\begin{practiceQuestion}{4-2-q14}{mc}
\begin{questionText}
Simplify $\frac{5}{6}t - \frac{1}{3}t + 4$.
\end{questionText}
\choiceA{$\frac{1}{2}t + 4$}
\choiceB{$\frac{4}{3}t + 4$}
\choiceC{$t + 4$}
\choiceD{$\frac{2}{3}t + 4$}
\correctAnswer{A}
\explanation{Rewrite $\frac{1}{3}$ as $\frac{2}{6}$. Then $\frac{5}{6}t - \frac{2}{6}t = \frac{3}{6}t = \frac{1}{2}t$. Result: $\frac{1}{2}t + 4$.}
\end{practiceQuestion}

\begin{practiceQuestion}{4-2-q15}{mc}
\begin{questionText}
Simplify $3a + 4b - a + 2b$.
\end{questionText}
\choiceA{$4a + 6b$}
\choiceB{$2a + 6b$}
\choiceC{$2a + 2b$}
\choiceD{$4a + 2b$}
\correctAnswer{B}
\explanation{Combine $a$-terms: $3a - a = 2a$. Combine $b$-terms: $4b + 2b = 6b$. Result: $2a + 6b$.}
\end{practiceQuestion}

% --- Short Answer Questions (6) ---

\begin{practiceQuestion}{4-2-q16}{sa}
\begin{questionText}
Simplify $11c - 6 + 3c - 9$.
\end{questionText}
\correctAnswer{$14c - 15$}
\explanation{Combine variable terms: $11c + 3c = 14c$. Combine constants: $-6 - 9 = -15$. Result: $14c - 15$.}
\end{practiceQuestion}

\begin{practiceQuestion}{4-2-q17}{sa}
\begin{questionText}
Simplify $\frac{2}{3}m + \frac{1}{6}m - 5$.
\end{questionText}
\correctAnswer{$\frac{5}{6}m - 5$}
\explanation{Find a common denominator: $\frac{4}{6}m + \frac{1}{6}m = \frac{5}{6}m$. The constant stays: $\frac{5}{6}m - 5$.}
\end{practiceQuestion}

\begin{practiceQuestion}{4-2-q18}{sa}
\begin{questionText}
Simplify $-3p + 7q + 5p - 2q$.
\end{questionText}
\correctAnswer{$2p + 5q$}
\explanation{Combine $p$-terms: $-3p + 5p = 2p$. Combine $q$-terms: $7q - 2q = 5q$. Result: $2p + 5q$.}
\end{practiceQuestion}

\begin{practiceQuestion}{4-2-q19}{sa}
\begin{questionText}
Simplify $-7 + 4x - 2x + 10 - x$.
\end{questionText}
\correctAnswer{$x + 3$}
\explanation{Combine variable terms: $4x - 2x - x = x$. Combine constants: $-7 + 10 = 3$. Result: $x + 3$.}
\end{practiceQuestion}

\begin{practiceQuestion}{4-2-q20}{sa}
\begin{questionText}
A student says $5x + 3y$ simplifies to $8xy$. Explain the error.
\end{questionText}
\correctAnswer{$5x$ and $3y$ are not like terms; the expression is already simplified.}
\explanation{Like terms must have the same variable. $5x$ has variable $x$ and $3y$ has variable $y$, so they cannot be combined. The expression $5x + 3y$ is already in simplest form.}
\end{practiceQuestion}

\begin{practiceQuestion}{4-2-q21}{sa}
\begin{questionText}
Simplify $15 - 6a + 2a - 11 + 3a$.
\end{questionText}
\correctAnswer{$-a + 4$ or $4 - a$}
\explanation{Combine variable terms: $-6a + 2a + 3a = -a$. Combine constants: $15 - 11 = 4$. Result: $-a + 4$.}
\end{practiceQuestion}

% --- Graphical Multiple Choice Questions (3) ---

\begin{practiceQuestion}{4-2-q22}{gmc}
\begin{questionText}
The algebra tiles below represent an expression. Which is the simplified form?

\begin{center}
\begin{tikzpicture}[scale=0.55]
  % Positive x-tiles
  \fill[funBlue!40] (0,0) rectangle (2,1);
  \draw[thick] (0,0) rectangle (2,1);
  \node at (1,0.5) {$x$};
  \fill[funBlue!40] (2.3,0) rectangle (4.3,1);
  \draw[thick] (2.3,0) rectangle (4.3,1);
  \node at (3.3,0.5) {$x$};
  \fill[funBlue!40] (4.6,0) rectangle (6.6,1);
  \draw[thick] (4.6,0) rectangle (6.6,1);
  \node at (5.6,0.5) {$x$};
  \fill[funBlue!40] (6.9,0) rectangle (8.9,1);
  \draw[thick] (6.9,0) rectangle (8.9,1);
  \node at (7.9,0.5) {$x$};
  % Negative x-tiles
  \fill[funRed!40] (0,-1.5) rectangle (2,-0.5);
  \draw[thick] (0,-1.5) rectangle (2,-0.5);
  \node at (1,-1) {$-x$};
  % Positive unit tiles
  \fill[funGreen!40] (4,-1.5) rectangle (5,-0.5);
  \draw[thick] (4,-1.5) rectangle (5,-0.5);
  \node at (4.5,-1) {$1$};
  \fill[funGreen!40] (5.3,-1.5) rectangle (6.3,-0.5);
  \draw[thick] (5.3,-1.5) rectangle (6.3,-0.5);
  \node at (5.8,-1) {$1$};
  \fill[funGreen!40] (6.6,-1.5) rectangle (7.6,-0.5);
  \draw[thick] (6.6,-1.5) rectangle (7.6,-0.5);
  \node at (7.1,-1) {$1$};
  % Negative unit tiles
  \fill[funRed!40] (8.2,-1.5) rectangle (9.2,-0.5);
  \draw[thick] (8.2,-1.5) rectangle (9.2,-0.5);
  \node at (8.7,-1) {$-1$};
\end{tikzpicture}
\end{center}
\end{questionText}
\choiceA{$3x + 2$}
\choiceB{$5x + 2$}
\choiceC{$3x + 4$}
\choiceD{$3x - 2$}
\correctAnswer{A}
\explanation{There are $4$ positive $x$-tiles and $1$ negative $x$-tile: $4x - x = 3x$. There are $3$ positive unit tiles and $1$ negative unit tile: $3 - 1 = 2$. Simplified: $3x + 2$.}
\end{practiceQuestion}

\begin{practiceQuestion}{4-2-q23}{gmc}
\begin{questionText}
The sides of the shape below are labeled with expressions. What is the simplified perimeter?

\begin{center}
\begin{tikzpicture}[scale=0.85]
  \draw[thick, fill=funBlue!10] (0,0) rectangle (4,2.5);
  \node[below] at (2,0) {$4x + 1$};
  \node[above] at (2,2.5) {$4x + 1$};
  \node[left] at (0,1.25) {$x + 3$};
  \node[right] at (4,1.25) {$x + 3$};
\end{tikzpicture}
\end{center}
\end{questionText}
\choiceA{$10x + 8$}
\choiceB{$5x + 4$}
\choiceC{$8x + 4$}
\choiceD{$10x + 4$}
\correctAnswer{A}
\explanation{Add all four sides: $(4x + 1) + (x + 3) + (4x + 1) + (x + 3) = 10x + 8$.}
\end{practiceQuestion}

\begin{practiceQuestion}{4-2-q24}{gmc}
\begin{questionText}
A student made the work shown below. In which step did the student make an error?

\begin{center}
\begin{tabular}{|l|l|}
\hline
\textbf{Step} & \textbf{Work} \\
\hline
Original & $6x + 3 - 2x - 5$ \\
\hline
Step 1 & $(6x + 2x) + (3 - 5)$ \\
\hline
Step 2 & $8x + (-2)$ \\
\hline
Step 3 & $8x - 2$ \\
\hline
\end{tabular}
\end{center}
\end{questionText}
\choiceA{Step 1: the student grouped the terms incorrectly}
\choiceB{Step 2: $6x + 2x$ should be $8x^2$}
\choiceC{Step 1: it should be $(6x - 2x)$, not $(6x + 2x)$}
\choiceD{Step 3: the answer should be $8x + 2$}
\correctAnswer{C}
\explanation{In the original expression the $x$-terms are $6x$ and $-2x$. The student wrote $6x + 2x$ instead of $6x - 2x$. The correct simplification is $4x - 2$.}
\end{practiceQuestion}

% --- Graphical Short Answer Questions (3) ---

\begin{practiceQuestion}{4-2-q25}{gsa}
\begin{questionText}
The diagram shows a triangle with the side lengths given. Write the simplified expression for the perimeter.

\begin{center}
\begin{tikzpicture}[scale=0.9]
  \coordinate (A) at (0,0);
  \coordinate (B) at (5,0);
  \coordinate (C) at (2.5,3);
  \draw[thick, fill=funGreen!10] (A) -- (B) -- (C) -- cycle;
  \node[below] at (2.5,0) {$5x - 2$};
  \node[right] at (3.9,1.65) {$2x + 7$};
  \node[left] at (1,1.65) {$3x + 1$};
\end{tikzpicture}
\end{center}
\end{questionText}
\correctAnswer{$10x + 6$}
\explanation{Add all three sides: $(5x - 2) + (2x + 7) + (3x + 1)$. Combine variable terms: $5x + 2x + 3x = 10x$. Combine constants: $-2 + 7 + 1 = 6$. Perimeter: $10x + 6$.}
\end{practiceQuestion}

\begin{practiceQuestion}{4-2-q26}{gsa}
\begin{questionText}
The two groups of algebra tiles below are placed together. Write the simplified expression they represent.

\begin{center}
\begin{tikzpicture}[scale=0.55]
  % Group 1
  \node[font=\bfseries] at (4, 1.5) {Group 1};
  \fill[funBlue!40] (0,0) rectangle (2,1);
  \draw[thick] (0,0) rectangle (2,1);
  \node at (1,0.5) {$x$};
  \fill[funBlue!40] (2.3,0) rectangle (4.3,1);
  \draw[thick] (2.3,0) rectangle (4.3,1);
  \node at (3.3,0.5) {$x$};
  \fill[funBlue!40] (4.6,0) rectangle (6.6,1);
  \draw[thick] (4.6,0) rectangle (6.6,1);
  \node at (5.6,0.5) {$x$};
  \fill[funGreen!40] (7,0) rectangle (8,1);
  \draw[thick] (7,0) rectangle (8,1);
  \node at (7.5,0.5) {$1$};

  % Group 2
  \node[font=\bfseries] at (4, -1) {Group 2};
  \fill[funRed!40] (0,-2.5) rectangle (2,-1.5);
  \draw[thick] (0,-2.5) rectangle (2,-1.5);
  \node at (1,-2) {$-x$};
  \fill[funRed!40] (2.3,-2.5) rectangle (4.3,-1.5);
  \draw[thick] (2.3,-2.5) rectangle (4.3,-1.5);
  \node at (3.3,-2) {$-x$};
  \fill[funGreen!40] (5,-2.5) rectangle (6,-1.5);
  \draw[thick] (5,-2.5) rectangle (6,-1.5);
  \node at (5.5,-2) {$1$};
  \fill[funGreen!40] (6.3,-2.5) rectangle (7.3,-1.5);
  \draw[thick] (6.3,-2.5) rectangle (7.3,-1.5);
  \node at (6.8,-2) {$1$};
  \fill[funGreen!40] (7.6,-2.5) rectangle (8.6,-1.5);
  \draw[thick] (7.6,-2.5) rectangle (8.6,-1.5);
  \node at (8.1,-2) {$1$};
\end{tikzpicture}
\end{center}
\end{questionText}
\correctAnswer{$x + 4$}
\explanation{Group 1: $3x + 1$. Group 2: $-2x + 3$. Combined: $(3x - 2x) + (1 + 3) = x + 4$.}
\end{practiceQuestion}

\begin{practiceQuestion}{4-2-q27}{gsa}
\begin{questionText}
An L-shaped figure is formed by removing a small rectangle from a larger rectangle. Write a simplified expression for the perimeter of the L-shape using the dimensions shown.

\begin{center}
\begin{tikzpicture}[scale=0.7]
  \draw[thick, fill=funBlue!10] (0,0) -- (6,0) -- (6,2) -- (3,2) -- (3,4) -- (0,4) -- cycle;
  \node[below] at (3,0) {$6x$};
  \node[right] at (6,1) {$2x + 1$};
  \node[above] at (4.5,2) {$3x$};
  \node[right] at (3,3) {$2x - 1$};
  \node[above] at (1.5,4) {$3x$};
  \node[left] at (0,2) {$4x$};
\end{tikzpicture}
\end{center}
\end{questionText}
\correctAnswer{$20x$}
\explanation{Add all sides: $6x + (2x + 1) + 3x + (2x - 1) + 3x + 4x = 20x + 0 = 20x$.}
\end{practiceQuestion}
